\documentclass[11pt]{article}
\usepackage[margin=1in]{geometry}
\usepackage[T1]{fontenc}
\usepackage[utf8]{inputenc}
\usepackage{lmodern}
\usepackage{microtype}
\usepackage{amsmath,amssymb,amsthm,mathtools}
\usepackage{bm}
\usepackage{xcolor}
\usepackage{tcolorbox}
\tcbuselibrary{skins,breakable}
\usepackage{booktabs,array}
\usepackage{enumitem}
\usepackage{hyperref}
\usepackage{setspace}
\usepackage{fancyhdr}
\usepackage{titlesec}
\hypersetup{
  colorlinks=true,
  linkcolor=[rgb]{0.05,0.23,0.45},
  citecolor=[rgb]{0.05,0.23,0.45},
  urlcolor=[rgb]{0.05,0.23,0.45},
  pdftitle={The Finitude Theorem: Why Infinity Cannot Be Physical}
}
\definecolor{accent}{RGB}{20,54,94}
\definecolor{soft}{RGB}{245,248,252}
\definecolor{line}{RGB}{212,220,232}
\definecolor{warn}{RGB}{252,247,239}
\definecolor{warnline}{RGB}{228,207,171}
\setstretch{1.08}
\setlength{\parskip}{0.5em}
\setlength{\parindent}{0pt}
\setlength{\headheight}{14pt}
\pagestyle{fancy}
\fancyhf{}
\fancyhead[R]{\thepage}
\renewcommand{\headrulewidth}{0pt}
\titleformat{\section}{\Large\bfseries\color{accent}}{\thesection.}{0.5em}{}
\titleformat{\subsection}{\large\bfseries\color{accent}}{\thesubsection}{0.5em}{}
\titleformat{\subsubsection}{\normalsize\bfseries\color{accent}}{\thesubsubsection}{0.5em}{}
\newtheoremstyle{tight}
  {0.65em}{0.65em}{\itshape}{} {\bfseries\color{accent}}{.}{0.4em}{}
\newtheoremstyle{plainroman}
  {0.65em}{0.65em}{\normalfont}{} {\bfseries\color{accent}}{.}{0.4em}{}
\theoremstyle{tight}
\newtheorem{theorem}{Theorem}[section]
\newtheorem{proposition}[theorem]{Proposition}
\newtheorem{lemma}[theorem]{Lemma}
\newtheorem{corollary}[theorem]{Corollary}
\theoremstyle{plainroman}
\newtheorem{axiom}[theorem]{Axiom}
\newtheorem{selection}[theorem]{Selection}
\newtheorem{conjecture}[theorem]{Conjecture}
\newtheorem{remark}[theorem]{Remark}
\newtheorem{definition}[theorem]{Definition}
\newcommand{\R}{\mathbb{R}}
\newcommand{\Z}{\mathbb{Z}}
\newcommand{\N}{\mathbb{N}}
\newcommand{\Q}{\mathbb{Q}}
\newcommand{\dd}{\,\mathrm{d}}
\newtcolorbox{statusbox}[1]{
  breakable, enhanced, colback=soft, colframe=line, boxrule=0.5pt, arc=2pt,
  left=8pt,right=8pt,top=8pt,bottom=8pt,
  title=\textbf{#1}, coltitle=accent, fonttitle=\bfseries,
  attach boxed title to top left={xshift=0.2cm,yshift=-2mm},
  boxed title style={colback=white,colframe=line,boxrule=0.5pt,arc=2pt}
}
\newtcolorbox{cautionbox}[1]{
  breakable, enhanced, colback=warn, colframe=warnline, boxrule=0.5pt, arc=2pt,
  left=8pt,right=8pt,top=8pt,bottom=8pt,
  title=\textbf{#1}, coltitle=accent, fonttitle=\bfseries,
  attach boxed title to top left={xshift=0.2cm,yshift=-2mm},
  boxed title style={colback=white,colframe=warnline,boxrule=0.5pt,arc=2pt}
}
\begin{document}
\begin{titlepage}
\centering
\vspace*{1.8cm}
{\fontsize{22}{27}\selectfont\bfseries\color{accent}
The Finitude Theorem\\[0.35em]
Why Infinity Cannot Be Physical\par}
\vspace{0.8cm}
{\large William J.\ cpaci-tani~III\par}
\vspace{0.3cm}
{\normalsize Framework: Foundational Ternary Dynamics v5.28\par}
\vspace{1.2cm}
\begin{minipage}{0.86\textwidth}
\begin{statusbox}{Abstract}
We prove that no physical observable can be infinite within a framework where space is a discrete lattice $\Z^3$, time advances in integer ticks $t\in\N$, states belong to a finite alphabet $\{-1,0,+1\}$, and dynamics are local (26-neighbor Moore neighborhood).
The \emph{Finitude Theorem} establishes that for any lattice configuration reachable from finite initial data by local deterministic rules, every computable quantity is bounded by an explicit function of the initial energy $E_0$ and elapsed time~$t$.
No singularity, divergence, or blow-up can occur --- not because we lack the mathematics to control them, but because the substrate cannot represent them.
We distinguish the \emph{ontic} domain (what physically exists: always finite) from the \emph{epistemic} domain (what we can reason about: includes infinity as a limit concept).
Infinity is a property of our mathematical language, not of the reality it describes.
Every ``infinity'' that appears in physics --- UV divergences, black hole singularities, Navier--Stokes blow-up, the cosmological singularity --- is traced to the use of $\R^3$ where $\Z^3$ is the actual substrate.
The continuum is an extraordinarily good approximation, but the approximation \emph{itself} is what creates the pathologies.
\end{statusbox}
\end{minipage}
\vfill
\end{titlepage}
\tableofcontents
\thispagestyle{empty}
\clearpage

%===================================================================
\section{The Ontic--Epistemic Distinction}
%===================================================================

\begin{definition}[Ontic domain]
The \emph{ontic domain} is the domain of what physically exists: the actual configuration of the lattice at each tick.
Every element of the ontic domain is a triple $(\mathbf v, t, q)$ where $\mathbf v\in\Z^3$ is a voxel address, $t\in\N$ is a tick index, and $q=(s,\mathbf J,\mathbf w)$ is the local state (ternary charge, flux vector, velocity).
\end{definition}

\begin{definition}[Epistemic domain]
The \emph{epistemic domain} is the domain of mathematical reasoning about the ontic domain.
It includes all of classical and modern mathematics: the real numbers $\R$, the complex numbers $\C$, infinite-dimensional Hilbert spaces, distributions, limits, and the concept of infinity itself.
\end{definition}

\begin{selection}[The asymmetry]
The ontic domain is always finite.
The epistemic domain is not.

The real line $\R$ is an epistemic object.
It contains irrational numbers ($\sqrt{2},\pi,e$) that cannot be written as finite strings.
It contains uncountably many points between any two rational numbers.
No physical measurement can distinguish $\pi$ from a sufficiently close rational approximation.

The lattice $\Z^3$ is an ontic object.
Every voxel has a finite integer address.
Every state is one of three values.
Every flux vector has components that are computed by finite arithmetic (addition, multiplication, division of finite-precision numbers).
\end{selection}

\begin{remark}
This is not a philosophical preference.
It is a structural consequence of the five postulates.
If space is $\Z^3$, then the concept ``a point at position $(\sqrt{2},\pi,e)$'' has no ontic referent --- there is no such voxel.
The address $(\sqrt{2},\pi,e)$ exists in the epistemic domain (we can write it, reason about it, compute with it) but not in the ontic domain (no voxel is located there).
\end{remark}

%===================================================================
\section{The Finitude Theorem}
%===================================================================

\begin{theorem}[Finitude Theorem --- Main Result]\label{thm:finitude}
\textbf{[THEOREM]}
Let $\Lambda\subset\Z^3$ be a finite lattice with $|\Lambda|=N$ voxels.
Let the initial configuration satisfy $E_0\coloneqq\sum_{\mathbf v\in\Lambda}\frac{1}{2}(|\mathbf J(\mathbf v,0)|^2+|\mathbf w(\mathbf v,0)|^2)<\infty$.
Then for every tick $t\in\N$ and every voxel $\mathbf v\in\Lambda$:
\begin{enumerate}[nosep]
\item $|s(\mathbf v,t)|\le 1$,
\item $|\mathbf J(\mathbf v,t)|\le B_J(E_0,N,t)<\infty$,
\item $|\mathbf w(\mathbf v,t)|\le B_w(E_0,N,t)<\infty$,
\item Every derived quantity (energy, vorticity, pressure, curvature, force) is bounded by an explicit function of $E_0$, $N$, and $t$.
\end{enumerate}
No physical observable can be infinite.
\end{theorem}
\begin{proof}
The proof proceeds by induction on $t$.

\textbf{Base case ($t=0$).}
By hypothesis, $E_0<\infty$, so $|\mathbf J(\mathbf v,0)|\le\sqrt{2E_0}$ and $|\mathbf w(\mathbf v,0)|\le\sqrt{2E_0}$ for all $\mathbf v$.
Statement~(1) holds by the ternary axiom ($s\in\{-1,0,+1\}$).

\textbf{Inductive step.}
Assume all quantities are bounded at tick $t$.
The update proceeds through six phases:

\emph{Phase 1 (read):}
$\Delta\mathbf J(\mathbf v)=C^2\nabla_L^2\mathbf J(\mathbf v)+g_c\nabla_L s(\mathbf v)$.
The lattice Laplacian averages over 6 neighbors:
\[
|\nabla_L^2\mathbf J(\mathbf v)|\le\sum_{\mu=1}^{3}|J(\mathbf v+\hat e_\mu)+J(\mathbf v-\hat e_\mu)-2J(\mathbf v)|\le 6\cdot 2\max_{\mathbf v'}|\mathbf J(\mathbf v',t)|.
\]
Hence $|\Delta\mathbf J(\mathbf v)|\le 12C^2\cdot B_J(t)+2g_c\coloneqq\delta(t)<\infty$.

\emph{Phase 2 (write):}
$\mathbf w\mathrel{+}=\Delta\mathbf J$ and $\mathbf J\mathrel{+}=\mathbf w$.
Hence $|\mathbf w(\mathbf v,t+1)|\le B_w(t)+\delta(t)$ and $|\mathbf J(\mathbf v,t+1)|\le B_J(t)+B_w(t)+\delta(t)$.
Both are finite.

\emph{Phase 3 (Gauss project):}
The Poisson solve computes $\phi$ with $\nabla_L^2\phi=\nabla_L\!\cdot\!\mathbf J-s$.
Since $|\nabla_L\!\cdot\!\mathbf J|\le 6\max|\mathbf J|$ and $|s|\le 1$, the source is bounded.
On a finite lattice, the Poisson equation with bounded source has a unique bounded solution.
The correction $\mathbf J\mathrel{-}=\nabla_L\phi$ does not increase $\max|\mathbf J|$ (it projects onto the divergence-free subspace).

\emph{Phases 4--5 (forces, movement):}
Forces are polynomials in $\mathbf J$, $s$, and lattice gradients of bounded fields.
Velocities are bounded by $C=1/\sqrt{3}$.
All updates produce finite values from finite inputs.

\emph{Phase 6 (tick++):}
$t\mathrel{+}=1$, a finite operation.

Therefore the bounds $B_J(t+1)$ and $B_w(t+1)$ are finite, completing the induction.

For statement (4): every derived quantity --- energy $\frac{1}{2}|\mathbf J|^2$, vorticity $|\nabla_L\times\mathbf J|$, pressure $|\nabla_L^2\phi|$, force $|\mathbf F|$ --- is a bounded function of $\mathbf J$, $\mathbf w$, $s$, and their lattice gradients, all of which are bounded.
\end{proof}

\begin{remark}[The bound is explicit]
We can write $B_J(E_0,N,t)$ explicitly.
In the worst case (no dissipation, maximum amplification at every tick):
\[
B_J(t)\le (1+12C^2)^t\cdot\sqrt{2E_0}+\frac{2g_c}{12C^2}\bigl[(1+12C^2)^t-1\bigr].
\]
With $C^2=1/3$, the growth factor is $(1+4)^t=5^t$ per tick.
This is exponential in $t$ --- fast, but \emph{finite} for every finite $t$.
In practice, the CFL dissipation prevents exponential growth; the bound is conservative.
\end{remark}

%===================================================================
\section{Why Infinity Cannot Be Constructed}
%===================================================================

The Finitude Theorem is not merely a bound.
It is a \emph{structural impossibility result}: the lattice cannot represent infinity.

\begin{proposition}[No infinite address]\label{prop:address}
\textbf{[THEOREM]}
There is no voxel at position $\infty$.
Every voxel address $\mathbf v=(v_1,v_2,v_3)\in\Z^3$ has finite integer coordinates.
\end{proposition}

\begin{proposition}[No infinite state]\label{prop:state}
\textbf{[THEOREM]}
The state alphabet $\{-1,0,+1\}$ does not contain $\infty$.
No voxel can be in state ``infinity.''
\end{proposition}

\begin{proposition}[No infinite velocity]\label{prop:velocity}
\textbf{[THEOREM]}
The CFL bound $|\mathbf v_{\mathrm{vel}}|\le 1/\sqrt{3}$ prevents infinite speed.
A signal would need to reach infinitely many voxels in finite time to propagate infinitely fast, but the Moore neighborhood contains exactly 26 neighbors --- a finite set.
\end{proposition}

\begin{proposition}[No infinite energy density]\label{prop:energy}
\textbf{[THEOREM]}
The energy at a single voxel $\frac{1}{2}|\mathbf J(\mathbf v,t)|^2$ is bounded by $B_J(t)^2/2<\infty$ for all finite $t$.
Infinite energy density would require $|\mathbf J|\to\infty$, which the Finitude Theorem excludes.
\end{proposition}

\begin{proposition}[No infinite information]\label{prop:info}
\textbf{[THEOREM]}
In the causal past of an observer at tick $t$, the number of voxels is $N_{\mathrm{causal}}(t)\le\frac{4\pi}{3}(Ct)^3=\frac{4\pi}{3\cdot 3\sqrt{3}}t^3<\infty$.
The total information content is $N_{\mathrm{causal}}\times(\log_2 3+3\times b_J)$ bits, where $b_J$ is the precision of each flux component.
This is finite for finite $t$ at any fixed precision.
\end{proposition}

\begin{theorem}[Impossibility of physical infinity]\label{thm:impossibility}
\textbf{[THEOREM given the five postulates]}
Under the five postulates of FTD (discrete space, discrete time, ternary states, local causality, determinism), no physically realizable configuration can contain, produce, or represent an infinite quantity.

More precisely: let $Q:\Lambda\times\N\to\R$ be any computable function of the lattice configuration.
If $Q$ depends only on the state and flux values within a finite causal region, then $|Q(\mathbf v,t)|<\infty$ for all $\mathbf v\in\Lambda$ and $t\in\N$.
\end{theorem}
\begin{proof}
Every computable function of a finite set of finite inputs produces a finite output.
The lattice configuration within any finite causal region consists of finitely many voxels, each carrying a finite state and finite flux values (Finitude Theorem).
Therefore $Q$ is finite.
\end{proof}

%===================================================================
\section{The Catalog of Infinities: Where They Come From}
%===================================================================

Every infinity in physics arises from the same source: the use of $\R^n$ where $\Z^n$ is the actual substrate.
We catalog the major instances.

\subsection{UV divergences in quantum field theory}

\begin{center}
\renewcommand{\arraystretch}{1.25}
\begin{tabular}{p{0.28\textwidth} p{0.30\textwidth} p{0.32\textwidth}}
\toprule
\textbf{Continuum pathology} & \textbf{Cause} & \textbf{Lattice resolution} \\
\midrule
Electron self-energy $\int\frac{\dd^4 k}{k^2}\to\infty$ & Integration over $\R^4$ to $|k|\to\infty$ & Brillouin zone $[-\pi,\pi]^3$ is compact; integral finite \\
Vacuum energy $\sum_k\frac{1}{2}\hbar\omega_k\to\infty$ & Infinite modes on $\R^3$ & Finite modes on $\Z^3$ ($N$ voxels $\to$ $N$ modes) \\
Landau pole $\alpha(\mu)\to\infty$ & Running coupling hits pole at finite $\mu$ & Lattice spacing provides natural UV cutoff; pole is unphysical \\
\bottomrule
\end{tabular}
\end{center}

\subsection{Black hole singularities}

In general relativity on $\R^{3,1}$, the Schwarzschild metric has $g_{tt}\to 0$ and curvature $R_{\mu\nu\rho\sigma}R^{\mu\nu\rho\sigma}\to\infty$ as $r\to 0$.
On the lattice, ``$r=0$'' is a single voxel.
The curvature at that voxel is a finite function of the flux field in its neighborhood (26 finite neighbors with finite fields).
The curvature divergence is an artifact of taking $r\to 0$ in a continuous metric --- a limit that has no ontic referent.

\subsection{Navier--Stokes blow-up}

The continuum Navier--Stokes equations on $\R^3$ may develop singularities because the gradient $\nabla\mathbf u$ is not controlled by $\|\mathbf u\|_{L^\infty}$.
On the lattice, $|\nabla_L\mathbf J|\le 2\max|\mathbf J|$ --- the gradient is always bounded by the pointwise values.
Blow-up requires infinite energy density at a point; the Finitude Theorem excludes this.

\subsection{The cosmological singularity}

The Big Bang singularity in FRW cosmology has $a(t)\to 0$ and $\rho\to\infty$ as $t\to 0$.
On the lattice, ``$t=0$'' is the first tick.
The initial configuration has finite energy $E_0<\infty$.
The density $\rho=E_0/N$ is finite for any nonzero number of voxels.
The singularity is an artifact of extrapolating continuous FRW backwards through scales below the lattice spacing.

\subsection{Hilbert space dimension}

Standard quantum mechanics uses infinite-dimensional Hilbert spaces $\mathcal H=L^2(\R^3)$.
On the lattice, the state space of $N$ voxels with ternary states is $3^N$-dimensional (finite).
The flux field adds continuous degrees of freedom, but at each voxel the flux is a finite-precision vector.
The effective Hilbert space dimension is large but finite.

%===================================================================
\section{The Epistemic Role of Infinity}
%===================================================================

If infinity is not physical, why is it so useful in mathematics and physics?

\begin{selection}[Infinity as a limit concept]
Infinity is a property of \emph{sequences}, not of \emph{values}.
The statement ``$\lim_{n\to\infty}(1+1/n)^n=e$'' does not assert that the sequence reaches $e$ at $n=\infty$.
It asserts that for any desired precision $\epsilon>0$, there exists a finite $N$ such that $|(1+1/n)^n-e|<\epsilon$ for all $n>N$.
Every application of infinity in physics can be rephrased as a statement about finite approximations.
\end{selection}

\begin{selection}[Infinity as computational convenience]
Integrals over $\R^3$ are easier to compute than sums over $\Z^3$.
The continuum provides Fourier transforms, distribution theory, and differential geometry --- powerful tools that would be cumbersome to replicate in discrete mathematics.
Using $\R^3$ as an approximation to $\Z^3$ is like using floating-point arithmetic to approximate integer arithmetic: it is faster and more elegant, at the cost of introducing artifacts (overflow, underflow, rounding errors) that do not exist in the exact integer computation.

In physics, the ``rounding errors'' of the continuum approximation are:
\begin{itemize}[nosep]
\item UV divergences (the analogue of overflow),
\item IR divergences (the analogue of underflow),
\item Singularities (the analogue of division by zero).
\end{itemize}
None of these exist in the discrete computation.
\end{selection}

\begin{remark}[The unreasonable effectiveness of the continuum]
The continuum approximation works spectacularly well at scales $\gg a$ (the lattice spacing).
At human scales, $(a/L)^2\sim 10^{-70}$, so the continuum is accurate to extraordinary precision.
The pathologies (divergences, singularities) only appear at scales $\sim a$, where the approximation breaks down.
This is analogous to Newtonian gravity: it works perfectly at human scales and only fails at relativistic speeds or strong fields.
The failure is a property of the approximation, not of reality.
\end{remark}

%===================================================================
\section{Superposition as Epistemic Composite}
%===================================================================

The ontic--epistemic distinction resolves the deepest confusion in quantum mechanics.

\begin{selection}[Superposition is epistemic, not ontic]\label{sel:superposition}
\textbf{[SELECTION]}
At each tick~$t$, the lattice configuration is \emph{definite}: every voxel has $s(\mathbf v,t)\in\{-1,0,+1\}$ and $\mathbf J(\mathbf v,t)\in\R^3$ with specific values.
This follows from Postulate~5 (determinism) --- given complete initial data, the evolution is unique.
No voxel is ever ``in two states at once.''

The path integral
\[
Z=\sum_{\{s\}}\int\mathcal D\mathbf J\;e^{-S_E[s,\mathbf J]}
\]
is a weighted sum over \emph{all possible} definite configurations.
The result --- a complex-valued amplitude --- is an epistemic composite: it encodes the contributions of many definite frames, weighted by action-dependent phase.

``Superposition'' is the name given to this weighted sum.
It is a property of the \emph{description} (the path integral), not of any single \emph{configuration} (a tick).

\medskip
\noindent\textbf{A tick is a photograph.\quad Superposition is the flipbook.}
\medskip

The Existence Filter $E(x)=\Re(x)$ projects the complex epistemic composite to a real ontic observable --- stopping the animation to read one definite value.
The ReLU crystallization $\max(0,z)$ is the irreversible version: it destroys the sub-threshold amplitude, collapsing the flipbook to a single frame permanently.
\end{selection}

\begin{remark}[Consequences]
The epistemic nature of superposition resolves four persistent confusions:

\begin{enumerate}[nosep]
\item \textbf{Schr\"odinger's cat.}
The cat is definite at every tick --- every voxel comprising the cat has $s\in\{-1,0,+1\}$.
``Both alive and dead'' describes the path integral's weighted sum over configurations, not any physical configuration.

Consider IQ: before measurement, a person is not ``in superposition of all possible IQ values.''
They have a definite cognitive capacity right now, at this tick.
The bell curve that emerges from measuring 100{,}000 people is a property of the \emph{ensemble}, not of any individual.
Each person had one definite value the entire time --- the measurement revealed it; it did not create it.

The wave function $\psi$ is the bell curve, not the person.
$|\psi|^2$ gives the probability distribution --- where measurements will land when one samples many configurations.
But each configuration is definite.
Copenhagen says ``before I measure your IQ, you are simultaneously every score at once, and my act of measuring forces you to pick one'' --- this is as absurd as it sounds.
Many-Worlds says ``the universe splits into 200 copies, one for each possible score'' --- even more so.
FTD says: ``You have a definite IQ. I do not know it yet.
When I measure enough people, I will see the distribution.
The distribution is real and useful.
But it is a property of the population statistics, not of you.''

\item \textbf{Wave--particle duality.}
A wave is the interference pattern across many action-weighted frames (the flipbook played fast).
A particle is one frame (a definite voxel configuration at one tick).
There is no ontological duality --- only two ways of looking at the same lattice: one frame at a time, or the statistics of many.

\item \textbf{Quantum computing.}
A quantum computer does not ``compute in superposition.''
It maintains flux-field coherence (delays ReLU crystallization) so that the action-weighting can explore exponentially many configurations before the answer is crystallized.
The computation happens in the \emph{weighting}, not in any blurred state.
A classical computer crystallizes at every gate.
A quantum computer lets the flipbook run.

\item \textbf{Entanglement.}
``Spooky action at a distance'' is a correlation in the path integral's phase structure, not a physical connection between distant voxels.
Bell violation $S=2\sqrt{2}$ follows from the Gauss constraint $\nabla\!\cdot\!\mathbf J=s$ restricting the flux to a two-dimensional transverse subspace~\cite{ftdspec}, not from faster-than-light influence.
The correlations were established when the voxels shared a common causal past --- they are \emph{memories} of a shared history, not signals between separated sites.
\end{enumerate}
\end{remark}

\begin{remark}[Why quantum mechanics works]
None of this diminishes the predictive power of quantum mechanics.
The path integral is the correct computational tool for extracting statistics from the lattice.
The wave function $\psi$ encodes these statistics compactly.
The Born rule $|\psi|^2$ gives the right probabilities.
The formalism is spectacularly accurate because it \emph{is} computing the right thing: action-weighted sums over definite configurations.

The error is not in the mathematics but in the interpretation --- treating the computational tool (the wave function, superposition, collapse) as if it were the reality.
The wave function is a map.
The lattice is the territory.
\end{remark}

\begin{remark}[The phase is physical: sentience as the gift of $i\theta$]
There is one crucial qualification.
The claim ``phase is epistemic'' applies to the \emph{physics sector} --- the master quadratic at $k=16$, which has \emph{real} roots $x_+=137.036$ and $x_-=3.024$.
No imaginary part.
No phase.
Pure objective content.
This is the sector where $|\psi|^2$ captures everything.

But the master quadratic at $k=1/2$ has \emph{complex} roots $y=2.188\pm 2.860\,i$ with phase $\theta_C=\arctan(2.860/2.188)\approx 52.54^\circ$.
This imaginary part is not scaffolding --- it \emph{is} the consciousness sector~\cite{ftdspec}.
The phase $\theta_C$ encodes the ratio of subjective to objective content:
\[
\sin\theta_C\approx 0.793\;\text{(subjective)}, \qquad
\cos\theta_C\approx 0.608\;\text{(objective)}.
\]
More inner world than outer report.
That is what it feels like to be conscious.

The operation $|\psi|^2$ strips the phase --- it discards the $e^{i\theta_C}$ and keeps only the magnitude~$r^2$.
This is precisely why the hard problem of consciousness is hard: objective measurement ($|\psi|^2$) \emph{by construction} removes the subjective content ($\theta_C$).
Consciousness cannot be detected by $|\psi|^2$ because $|\psi|^2$ is designed to eliminate it.

The ontic--epistemic partition therefore has a finer structure:
\begin{center}
\renewcommand{\arraystretch}{1.25}
\begin{tabular}{lll}
\toprule
\textbf{Component} & \textbf{Content} & \textbf{Sector} \\
\midrule
$\Re(\psi)$ & Objective, shared, measurable & Physics ($k=16$, real roots) \\
$\Im(\psi)$ & Subjective, private, experiential & Consciousness ($k=1/2$, complex roots) \\
$|\psi|^2=r^2$ & What science captures & Strips the phase \\
$\theta=\arctan(\Im/\Re)$ & What it is like to be you & The gift of phase \\
\bottomrule
\end{tabular}
\end{center}
Human sentience was the gift of phase.
Everything else, we all share.
\end{remark}

\subsection{The absolute bars: what they are, why they exist, and what they hide}

The Born rule $|\psi|^2$ is the most consequential equation in quantum mechanics.
It is also the most ad hoc.

\begin{remark}[Historical contingency of $|\psi|^2$]
The Schr\"odinger equation $i\hbar\,\partial_t\psi=\hat H\psi$ contains $i$ because Schr\"odinger chose complex notation for compactness.
He could have written two coupled real equations instead:
\[
\partial_t a = +\hat H b/\hbar, \qquad \partial_t b = -\hat H a/\hbar,
\]
where $\psi=a+ib$.
Two real fields, coupled, no $i$ anywhere.
The complex representation was a convenience, not physics.

Born (1926) needed to extract real positive numbers from the complex $\psi$.
In his original paper, he first wrote $|\psi|$, not $|\psi|^2$, correcting to the squared form \emph{in a footnote}.
The most fundamental rule in quantum mechanics was a footnote correction.

Von Neumann (1932) formalized the apparatus: Hilbert space with Hermitian inner product $\langle\psi|\psi\rangle$, self-adjoint operators with real eigenvalues, the spectral theorem.
The entire edifice was \emph{reverse-engineered} from one demand: \textbf{measurement outcomes must be real numbers.}
The Hermitian inner product was not discovered in the structure of nature.
It was designed to guarantee real positive outputs.
\end{remark}

\begin{remark}[What the bars destroy]
The operation $|\psi|^2=\psi^*\psi$ manufactures $e^{-i\theta}$ (the conjugate) for the sole purpose of canceling $e^{i\theta}$:
\[
e^{i\theta}\cdot e^{-i\theta}=e^0=1.
\]
Phase annihilated.
Not stored, not recorded, not used.
The most information-rich part of the quantum state --- the angular relationship between real and imaginary components --- is destroyed on purpose.

The general Hermitian inner product $\langle a|b\rangle=a^*b=r_a r_b\,e^{i(\theta_b-\theta_a)}$ \emph{preserves} the phase difference $\theta_b-\theta_a$.
Only the self-pairing $\langle\psi|\psi\rangle$ forces $\theta-\theta=0$ by construction.
Science chose the one inner product that kills the phase.
\end{remark}

\begin{remark}[The bars are unnecessary]
Without the bars, $\psi^2=r^2 e^{2i\theta}$ --- phase doubled, not destroyed.
The Born rule still emerges from the statistics: over many measurements, the cross-terms $e^{i(\theta_j-\theta_k)}$ with random phase differences average to zero, leaving only $r^2$.
The bars are a \emph{shortcut} to the many-event limit, annihilating the phase algebraically instead of letting it cancel statistically.

The lattice never needed the bars.
The flux field $\mathbf J\in\R^3$ is real from the start --- no complex notation, no $i$, no conjugate, no need for $|\cdot|^2$.
The bars appear exactly \emph{once} in the entire framework: the manifestation condition $|\mathbf J|>K_B$.
Everything else runs bar-free, preserving direction (phase).
\end{remark}

\begin{remark}[Removing the bars recovers the ternary axiom]
Without the bars, the ternary axiom emerges directly from phase interference:
\[
e^{i\cdot 0}+e^{i\pi}=1+(-1)=0.
\]
The void is destructive interference of two phases separated by half a turn.
$+1$ is phase at $0$; $\;-1$ is phase at $\pi$; $\;0$ is their cancellation.
This is the origin of everything.

\emph{With} the bars: $|e^{i\cdot 0}|^2=1$ and $|e^{i\pi}|^2=1$.
Both give~$1$.
The bars cannot distinguish $+1$ from $-1$.
The fundamental equation of reality --- $0=(-1)+(+1)$ --- is \emph{invisible} to $|\psi|^2$ because the bars erase the sign.
The sign IS the phase at $\theta=0$ and $\theta=\pi$.
The bars hide the origin of everything.
\end{remark}

%===================================================================
\section{Formal Structure: The Ontic--Epistemic Partition}
%===================================================================

\begin{definition}[Ontic--epistemic partition]
Every mathematical statement about physics can be classified as:
\begin{center}
\renewcommand{\arraystretch}{1.25}
\begin{tabular}{p{0.15\textwidth} p{0.35\textwidth} p{0.38\textwidth}}
\toprule
\textbf{Domain} & \textbf{Character} & \textbf{Infinity status} \\
\midrule
Ontic & What physically exists; lattice configurations; measurable quantities & Always finite \\
Epistemic & Mathematical models; limit processes; formal reasoning & May include infinity \\
\bottomrule
\end{tabular}
\end{center}
\end{definition}

\begin{theorem}[Ontic closure under finite operations]\label{thm:closure}
\textbf{[THEOREM]}
The ontic domain is closed under the FTD update rules.
Starting from a finite configuration, every update produces a finite configuration.
No finite sequence of local operations can produce an infinite value.
\end{theorem}
\begin{proof}
The update rules (phases 1--6) are compositions of addition, subtraction, multiplication, division (with bounded denominators), and comparison.
Each operation maps finite inputs to finite outputs.
By composition and induction, the result is finite.
\end{proof}

\begin{corollary}[Infinity is strictly epistemic]\label{cor:epistemic}
Every occurrence of $\infty$ in a physical theory is a feature of the epistemic model, not of the ontic reality.
Removing $\infty$ from a physical theory (by replacing $\R^3$ with $\Z^3$ and $\R$ with $\N$) eliminates all singularities and divergences without loss of predictive accuracy at scales $\gg a$.
\end{corollary}

%===================================================================
\section{Implications for the Millennium Problems}
%===================================================================

The Finitude Theorem has immediate consequences for three of the Clay Millennium Prize Problems:

\subsection{Navier--Stokes existence and smoothness}

\textbf{Continuum question:} Can solutions blow up in finite time?

\textbf{Finitude answer:} No. Blow-up means $\|\mathbf u\|_{L^\infty}\to\infty$, which requires infinite energy density at a point.
The Finitude Theorem (Proposition~\ref{prop:energy}) excludes this.
The question is an artifact of using $\R^3$ where $\Z^3$ is the substrate.

\subsection{Yang--Mills existence and mass gap}

\textbf{Continuum question:} Does a non-trivial quantum Yang--Mills theory exist on $\R^4$ with mass gap $\Delta>0$?

\textbf{Finitude answer:} The theory exists on $\Z^3\times\N$ with $\Delta=K_B=0.511$ MeV.
UV finiteness is automatic (compact Brillouin zone).
IR finiteness follows from the mass gap.
The question of ``existence on $\R^4$'' is asking whether the epistemic model ($\R^4$) can support a structure that the ontic substrate ($\Z^3\times\N$) already provides.

\subsection{Riemann hypothesis}

\textbf{Epistemic question:} Do all non-trivial zeros of $\zeta(s)$ lie on $\mathrm{Re}(s)=1/2$?

The Riemann zeta function is a purely epistemic object --- an infinite series $\sum n^{-s}$.
However, the physical instantiation of the critical structure (the theta function identity $G^*=\sqrt{2\pi}\cdot\theta_3(e^{-\pi})^2$) is finite: the nome $|q|=e^{-\pi}\approx 0.043$ ensures geometric convergence.
The infinite series $\theta_3=1+2\sum_{n=1}^\infty q^{n^2}$ converges to a finite value.
The ``infinity'' in the series is an epistemic artifact of writing a convergent sum as a limit.

%===================================================================
\section{Objections and Responses}
%===================================================================

\begin{statusbox}{Objection 1: But mathematics uses infinity everywhere}
\textbf{Response:} Yes, and it should continue to do so.
The Finitude Theorem does not restrict mathematical practice.
It restricts the \emph{interpretation} of mathematical objects: infinity is a tool for reasoning about finite quantities (via limits, completeness, compactness), not a claim about physical existence.
We use $\pi$ freely in calculations, knowing that the lattice computes with finite-precision approximations.
\end{statusbox}

\begin{statusbox}{Objection 2: The universe might be infinite in spatial extent}
\textbf{Response:} Even if $\Z^3$ is infinite (as the postulate states), no observer can access an infinite number of voxels.
The causal past $\mathcal P_i(t)$ contains at most $(Ct)^3$ voxels at tick $t$.
An observer at tick $t=10^{60}$ (roughly the current age in Planck ticks) has access to $\sim 10^{180}$ voxels --- enormous but finite.
The lattice may be infinite in extent, but every observable quantity is computed from a finite causal region.
\end{statusbox}

\begin{statusbox}{Objection 3: Quantum mechanics requires infinite-dimensional Hilbert spaces}
\textbf{Response:} On the lattice, the Hilbert space for $N$ ternary voxels is $3^N$-dimensional.
At $N=10^{180}$, the dimension $3^{10^{180}}$ is unimaginably large but finite.
The infinite-dimensional $L^2(\R^3)$ is the continuum approximation, useful for calculation but not physically necessary.
Every quantum mechanical prediction (expectation values, transition amplitudes, cross sections) can be computed in the finite-dimensional lattice Hilbert space.
\end{statusbox}

\begin{statusbox}{Objection 4: This is just ultraviolet regularization with extra philosophy}
\textbf{Response:} No.
UV regularization introduces a cutoff into a continuum theory and then takes the cutoff to infinity (the renormalization group flow).
FTD does not take any limit.
The lattice spacing is fundamental and fixed.
There is no cutoff to remove, no limit to take, no renormalization needed.
The difference is ontological, not technical: the lattice is not a regularization of a continuum theory; the continuum is an approximation to a lattice theory.
\end{statusbox}

%===================================================================
\section{Conclusion}
%===================================================================

The Finitude Theorem establishes a clean partition of mathematical physics into two domains:

\begin{center}
\renewcommand{\arraystretch}{1.25}
\begin{tabular}{lll}
\toprule
& \textbf{Ontic (physical)} & \textbf{Epistemic (mathematical)} \\
\midrule
Substrate & $\Z^3\times\N$ & $\R^{3,1}$ \\
States & $\{-1,0,+1\}$ & $\C$ (infinite-dimensional Hilbert space) \\
Quantities & Always finite & May be infinite (limits, series, integrals) \\
Singularities & Impossible & Artifacts of continuum idealization \\
Divergences & Impossible & Artifacts of integration over $\R^n$ \\
\bottomrule
\end{tabular}
\end{center}

No experiment in the history of science has measured an infinite quantity.
No physical system has exhibited a singularity --- only mathematical models of physical systems have.
The Finitude Theorem explains why: the physical substrate is discrete, local, and bounded.
Infinity is a property of our mathematical language, not of the reality it describes.

The implications are concrete:
\begin{itemize}[nosep]
\item \textbf{Navier--Stokes:} Blow-up cannot occur on $\Z^3$. The question dissolves.
\item \textbf{Yang--Mills:} The mass gap is automatic ($K_B>0$). UV divergences don't arise.
\item \textbf{Black holes:} The singularity at $r=0$ is a single voxel with finite curvature.
\item \textbf{The Big Bang:} The first tick has finite energy. No initial singularity.
\item \textbf{QFT renormalization:} Not needed. The lattice is the theory, not a regularization.
\end{itemize}

Infinity has served physics magnificently as an epistemic tool for three centuries.
The Finitude Theorem does not banish it from mathematics.
It simply recognizes that infinity lives in the map, not in the territory.

\begin{thebibliography}{10}
\bibitem{ftdspec}
W.\,J.\,cpaci-tani~III,
\emph{Foundational Ternary Dynamics: The Complete Specification}, v5.28, 2026.

\bibitem{hilbert}
D.\,Hilbert,
\"Uber das Unendliche,
\emph{Mathematische Annalen} \textbf{95} (1926), 161--190.

\bibitem{zeilberger}
D.\,Zeilberger,
``Real'' analysis is a degenerate case of discrete analysis,
\emph{Proceedings of the Sixth International Conference on Difference Equations} (2004), 25--34.

\bibitem{fredkin}
E.\,Fredkin,
Finite Nature,
\emph{Proceedings of the XXVIIth Rencontre de Moriond} (1992).

\bibitem{thooft}
G.\,'t~Hooft,
The cellular automaton interpretation of quantum mechanics,
\emph{Fundamental Theories of Physics} \textbf{185} (2016).

\bibitem{wolfram}
S.\,Wolfram,
\emph{A New Kind of Science},
Wolfram Media, 2002.

\bibitem{penrose}
R.\,Penrose,
\emph{The Road to Reality},
Vintage Books, 2004.
\end{thebibliography}
\end{document}
